\section{Experimental Method}
\label{sec:evaluation}

Published Siderust v0.16.0 is evaluated through independent numerical comparisons,
executable controls and paired arithmetic measurements. The evaluated package
corresponds to release revision \texttt{3e0069d3}. A committed Cargo lockfile
pins the published crate versions used by the accompanying executable examples,
which supply every printed listing. Reference methods, test inputs and tolerances
are described below. Original raw measurements, independent reference
fixtures and reproduction scripts are preserved in a separate, version-pinned
supplement branch rather than in the publication-source tree.
The full experimental execution infrastructure is not required for building the manuscript.

\subsection{Reference transformations and ephemerides}

We compare time and frame operations with ERFA 2.0.1 through PyERFA 2.0.1.5,
and kernel-backed positions with CSPICE N0067 through SpiceyPy 6.0.0.
Earth rotation angle, IAU 2006 GMST, mean obliquity and the IAU 2006/2000A
bias--precession--nutation matrix are evaluated at 1,000 deterministic Julian
dates $2451545-36525+73.125i$, $0\leq i<1000$. These isolated rotation tests
supply equal numerical TT and UT1 dates as controlled arguments. Ten UTC dates
from 1972 to 2025 exercise UTC--TAI and UTC--TT around selected leap-second
changes; full-JD comparisons use a $50\,\mu$s tolerance.
A separate pinned feature graph enables runtime time-data loading for a compact
IERS extension: 13 civil UTC labels include the 1972, 2015 and 2016 inserted
seconds, and six epochs compare EOP-backed UT1, Earth rotation and GMST.
Independent ERFA two-part UTC quasi-JD references are converted to continuous
TAI, TT and UT1 before differencing. The native time interface retains split
seconds; scalar UTC JD is rejected during an inserted second.
IERS daily samples also test polar-motion values and leap-aware interpolation
of UT1--TAI across the 2016 step. The evaluation used fixed IERS inputs, documented coverage limits and
specified acceptance tolerances.

For ephemerides, we compare Earth heliocentric and Moon geocentric positions
with the same DE440s kernel \cite{park2021de440,naifspk}. The sample contains
1,000 eight-day-spaced TT epochs beginning at J2000 and one-second interior
points near both coverage endpoints. CSPICE converts TT to ET using the NAIF
leap-second kernel. ERFA's \texttt{ecm06} aligns the reference ICRF vectors
with Siderust's IAU 2006 mean-ecliptic J2000 convention; SPICE's
\texttt{ECLIPJ2000} uses a different obliquity. Exact endpoints and one-day
out-of-coverage requests are tested separately.

Acceptance bounds are $2\times10^{-11}$ rad for Earth rotation and GMST,
$10^{-14}$ rad for obliquity, $2\times10^{-12}$ per matrix element,
$2\times10^{-12}$ au for Earth and $2\times10^{-4}$ km for the Moon.
They measure implementation agreement with the specified references.
The analytical VSOP87A and ELP2000-82B outputs are compared separately with
DE440s as backend-output differences. VSOP87A retains its native dynamical
ecliptic and equinox of J2000; ELP2000-82B retains its dynamical ecliptic J2000
coordinates after the implemented Laskar rotation
\cite{bretagnon1988vsop87,chapront1988elp85}. Both are returned with the
\texttt{EclipticMeanJ2000} marker in v0.16.0, while the SPK path uses the
IAU 2006 realization above. The recorded comparison therefore includes
frame-realization effects as well as orbital-model differences.

For ephemeris positions, the residual is the Euclidean norm
\begin{equation}
 \Delta r=\left\|\boldsymbol r_{\mathrm{Siderust}}
                  -\boldsymbol r_{\mathrm{reference}}\right\|_2.
\end{equation}
Matrix residuals are maximum absolute element differences. Altitude and time
residuals are absolute scalar differences; azimuth differences are wrapped
to the shorter angular separation.

\subsection{Surface geometry and integrated observations}

The terrestrial application is evaluated on its complete 289-epoch grid.
Thirty lunar direction cases combine five epochs in TT spanning 1900--2050,
two sites and three ICRS directions. A fixed synthetic finite target gives
ten distinct epoch/site/target geometries, each executed three times. Independent
CSPICE \texttt{IAU\_MOON} rotations and separately constructed local bases
provide references for angles, range and Cartesian components. Four Mars cases
exercise ellipsoid intersections and surface-normal verticals under both
planetographic and planetocentric latitude conventions. The reference surface
is an oblate ellipsoid of revolution with equatorial radius
$a=3\,396\,190$ m and polar radius $b=3\,376\,200$ m, matching
\texttt{BODY499\_RADII} in NAIF \texttt{pck00011.tpc} \cite{naifpck00011}.
The Earth-from-Moon application additionally combines
CSPICE DE440s geometry, lunar orientation and observer displacement in one
reference calculation.

Lunar acceptance bounds are $10^{-10}$ per rotation-matrix element,
$2\times10^{-8}$ degrees for angles, $10^{-3}$ m for range and $0.1$ m for
Cartesian components. Surface position, Cartesian-component and basis residuals
use the maximum absolute component difference; range uses the absolute scalar
difference, and angular rows take the maximum altitude/wrapped-azimuth residual.
Mars geometry bounds are $10^{-8}$ m for positions and
$10^{-14}$ for basis components. A separate cardinal test checks the Cartesian
north/east/up convention. The numerical comparisons used the same stated geometry and model
conventions for each case.

An additional analytic boundary suite isolates local geometry from orientation.
Its 48 lunar projections cover six sites, including exact and near poles,
longitudes at or near $\pm180^\circ$, heights from $-100$ to $10^4$ m,
four cardinal horizon directions, altitudes $\pm10^{-6}$ degrees, and
zenith/nadir. Finite slant ranges are $10^3$, $10^6$ and $3.8\times10^8$ m.
Targets are constructed in the body-fixed frame from an independent analytic
spherical basis. Twelve vertical cases expect the documented zero azimuth
convention, since physical azimuth is undefined. A further 32 Mars cases
combine eight sites, both latitude conventions and two heights on this
oblate ellipsoid. Independent radial intersections and surface-normal formulas
supply positions and bases. Bounds allow floating-point rounding: $2\times10^{-6}$ m
for lunar components/ranges, $2\times10^{-8}$ degrees for defined angles,
$10^{-8}$ m for Mars positions and $10^{-14}$ for basis components.
These tests isolate reusable geometry; the earlier CSPICE cases separately
exercise time-dependent lunar orientation.

\subsection{Cost of dimensional interfaces}

We compare kilometre-to-metre conversion and distance divided by a three-second
duration with matched scalar \texttt{f64} arithmetic. Both implementations use
the same inputs and outputs, with an equality check on 1,000 inputs before
timing. A single Rust 1.93.1 release binary uses default optimization and
\texttt{black\_box} around inputs and outputs. Five warm-up rounds precede
30 paired rounds of $10^6$ operations, alternating measurement order.
We report median cost and sample standard deviation across rounds, including
loop and measurement overhead in both implementations. The recorded host is an
Intel Core Ultra 9 185H running Linux. These kernels characterize the cost of
two common quantity operations on that host.

Typed/scalar cost ratios are formed within each paired round. Their median and
empirical p5--p95 interval describe the observed variability, without assuming
independent typed and scalar samples or interpreting quantiles as confidence
intervals. The kernels include scalar arithmetic, quantity construction and
conversion, loop and measurement overhead; they exclude ephemeris evaluation,
I/O and observer transformations.
